\فصل{مفاهیم اولیه}

این فصل مفاهیمی را توضیح می‌دهد که برای دنبال‌کردن مسئله و آزمایش‌های پایان‌نامه ضروری‌اند. بحث بر چهار پرسش متمرکز است: بازشناس گفتار چه ورودی و خروجی‌ای دارد و چگونه ارزیابی می‌شود؛ چرا گفتار فارسی در کانال‌های تلفنی با داده آموزشی عمومی ناهمخوان است؛ آموزش چندشرایطی و بهسازی دوجریانه چگونه به این ناهمخوانی پاسخ می‌دهند؛ و چگونه می‌توان صوت بلند بدون رونویسی را به داده آموزشی تبدیل کرد. جزئیات مطالعات پیشین در فصل بعد و انتخاب‌های اجرایی، مانند نمایه‌های تخریب، معماری دقیق و برنامه آموزش، در فصل راهکار پیشنهادی می‌آیند.

دو مسئله اصلی این پایان‌نامه به هم مرتبط‌اند. از یک سو، داده‌های برچسب‌خورده فارسی از نظر حجم، نوع گفتار و شرایط ضبط محدود هستند؛ از سوی دیگر، کانال‌های تلفنی و \lr{VoIP} ویژگی‌های آکوستیکی گفتار را نسبت به شرایط معمول تغییر می‌دهند. برای مواجهه با این دو مسئله، مسیر داده‌محور پژوهش بر گسترش و متنوع‌سازی داده‌های آموزشی، از جمله با استفاده از گفتار شبه‌برچسب‌خورده، تمرکز دارد و در مسیر مدل‌محور نیز تلاش می‌شود با ترکیب نمای اصلی و نمای بهسازی‌شده گفتار، ضمن کاهش اثر تخریب‌های ناشی از کانال مخابراتی، اطلاعات مفید هر دو بازنمایی حفظ شود. ساختار این فصل نیز بر همین سه محور داده، دامنه و بازنمایی شکل گرفته است.

\قسمت{بازشناسی گفتار و معیارهای ارزیابی}

\زیرقسمت{بازشناسی انتهابه‌انتها و ویسپر}

بازشناسی خودکار گفتار\پاورقی{Automatic Speech Recognition (ASR)} وظیفه تبدیل سیگنال صوتی به متن را بر عهده دارد. در این مسئله، مدل تلاش می‌کند برای ورودی صوتی $X$، محتمل‌ترین دنباله متنی $Y$ را پیدا کند. در روش‌های انتهابه‌انتها، بخش‌هایی که در سامانه‌های کلاسیک به‌صورت جداگانه طراحی می‌شدند، در قالب یک مدل واحد و قابل‌آموزش ترکیب می‌شوند. تفاوت اصلی میان خانواده‌های مختلف این مدل‌ها نیز به نحوه هم‌ترازی صوت و متن و تولید دنباله خروجی برمی‌گردد \مرجع{li2022recent}.

ورودی صوتی پیش از ورود به بسیاری از بازشناس‌ها به بازنمایی زمان-فرکانس تبدیل می‌شود. تبدیل فوریه کوتاه‌مدت\پاورقی{Short-Time Fourier Transform (STFT)} سیگنال را در پنجره‌های کوتاه تحلیل می‌کند و بانک فیلتر مل انرژی فرکانسی را در مقیاسی فشرده‌تر نمایش می‌دهد. لگاریتم این انرژی‌ها ویژگی لاگ-مل را می‌سازد؛ بازنمایی‌ای که تغییرات زمانی و طیفی گفتار را بدون نگه‌داری شکل موج خام در اختیار مدل قرار می‌دهد.

ویسپر\پاورقی{Whisper} یک مدل رمزگذار-رمزگشای ترنسفورمری است که با نظارت ضعیف روی داده‌های چندزبانه و چندوظیفه‌ای آموزش دیده است. در این مدل، صوت تک‌کاناله $16$ کیلوهرتزی ابتدا به ویژگی‌های لاگ-مل تبدیل می‌شود، سپس رمزگذار آن را به بازنمایی‌های نهان نگاشت می‌کند و رمزگشا متن را به‌صورت خودبازگشتی تولید می‌کند \مرجع{radford2023robust}. پیش‌آموزش گسترده ویسپر، آن را به نقطه شروع مناسبی برای زبان‌های کم‌داده تبدیل کرده است، اما مدل ازپیش‌آموزش‌دیده معمولاً برای استفاده مستقیم در این زبان‌ها و دامنه‌های خاص دقت کافی ندارد. به همین دلیل، ریزتنظیم و سازگارکردن مدل با زبان و شرایط آکوستیکی هدف همچنان ضروری است؛ موضوعی که مطالعات انجام‌شده روی زبان‌های کم‌منبع و فارسی نیز به آن اشاره کرده‌اند \مرجع{liu2024whisper,rahimipour2024persian,namdarzadeh2026sps}.

\زیرقسمت{نرخ خطای واژه و نویسه}

نرخ خطای واژه\پاورقی{Word Error Rate (WER)} معیار اصلی ارزیابی بازشناسی گفتار است. در محاسبه این معیار پس از هم‌ترازکردن خروجی مدل با متن مرجع، تعداد جانشینی‌ها $S$، حذف‌ها $D$ و درج‌ها $I$ بر تعداد واژه‌های مرجع $N$ تقسیم می‌شود:

\begin{equation}
\mathrm{WER}=\frac{S+D+I}{N}.
\end{equation}

نرخ خطای نویسه\پاورقی{Character Error Rate (CER)} همان فاصله ویرایشی را در سطح نویسه محاسبه می‌کند. اگر $S_c$، $D_c$ و $I_c$ خطاهای نویسه‌ای و $N_c$ تعداد نویسه‌های مرجع باشند، آنگاه

\begin{equation}
\mathrm{CER}=\frac{S_c+D_c+I_c}{N_c}.
\end{equation}

گزارش هم‌زمان این دو معیار برای فارسی اهمیت دارد. تفاوت در فاصله، نیم‌فاصله یا اتصال یک پسوند می‌تواند یک واژه را در WER کاملاً نادرست کند، درحالی‌که بیشتر نویسه‌های آن درست‌اند. در مقابل، درج‌های طولانی یا تکرارشونده ممکن است رفتار دو معیار را به شکل متفاوتی تغییر دهند؛ WER و CER در حضور درج فراوان حتی از $100$ درصد نیز فراتر می‌روند.

مقایسه معتبر به نرمال‌سازی یکسان متن مرجع و فرضیه نیاز دارد. یکسان‌سازی «ي/ی» و «ك/ک»، قواعد فاصله و نیم‌فاصله، اعداد و نشانه‌گذاری مستقیماً بر شمار خطا اثر می‌گذارد. ازاین‌رو، نرمال‌ساز بخشی از پروتکل ارزیابی است و نباید برای هر مدل یا پس از مشاهده خروجی آن تغییر کند.


\قسمت{ناهمخوانی دامنه در گفتار فارسی و مخابراتی}

\زیرقسمت{کم‌منبعی فارسی و تنوع نوشتاری}

کم‌منبع‌بودن فارسی به معنای نبود داده نیست، بلکه به محدودیت داده برچسب‌خورده و قابل دسترس در مقایسه با زبان‌های پرداده اشاره دارد. پیکره‌هایی مانند فارس‌دَت زیرساخت پژوهش فارسی را گسترش داده‌اند، اما بخش مهمی از داده‌های قدیمی شامل گفتار خوانده‌شده، واژگان محدود یا ضبط کنترل‌شده است \مرجع{bijankhan1994farsdat,sameti2008nevisa,sameti2011large}. مجموعه‌های عمومی جدیدتر دسترسی و تنوع را افزایش داده‌اند، ولی همه سبک‌های گفتار بداهه و کانال‌های کاربردی را پوشش نمی‌دهند \مرجع{commonvoice2024,conneau2022fleurs}.

تفاوت‌های نوشتاری فارسی هم در متن مرجع و هم در خروجی مدل دیده می‌شوند و اگر نرمال‌سازی یکسانی روی آن‌ها اعمال نشود، می‌توانند مستقیماً بر معیارهای ارزیابی اثر بگذارند. برای نمونه، برخی حروف عربی و فارسی با ظاهر مشابه، کدهای یونی‌کد متفاوتی دارند؛ پسوندها و پیشوندهایی مانند «ها» و «می» ممکن است با فاصله، نیم‌فاصله یا به‌صورت پیوسته نوشته شوند؛ و برخی واژه‌های محاوره‌ای نیز بیش از یک صورت نوشتاری دارند. در نتیجه، بخشی از اختلاف میان خروجی مدل و متن مرجع ممکن است ناشی از تفاوت‌های صرفاً نوشتاری باشد، نه خطای واقعی در بازشناسی گفتار. نرمال‌سازی متن با یکسان‌سازی این موارد، اثر اختلاف‌های صوری را کاهش می‌دهد و ارزیابی را بیشتر به خطاهای واقعی بازشناسی معطوف می‌کند. این مسئله از ناهمخوانی آکوستیکی مستقل است و باید در مرحله آماده‌سازی متن و محاسبه معیارها به‌صورت جداگانه در نظر گرفته شود. جدول~\ref{tab:text-normalization} نمونه‌هایی از تفاوت‌های نوشتاری و صورت نرمال‌شده آن‌ها را نشان می‌دهد.

\begin{table}[htbp]
\centering
\caption{نمونه‌هایی از نرمال‌سازی متن فارسی پیش از ارزیابی}
\label{tab:text-normalization}
\begingroup
\scriptsize
\setlength{\tabcolsep}{4pt}
\renewcommand{\arraystretch}{1.35}
\begin{tabular}{|p{0.22\textwidth}|p{0.38\textwidth}|p{0.28\textwidth}|}
\hline
\textbf{نوع تفاوت} & \textbf{صورت‌های ممکن} & \textbf{صورت نرمال‌شده} \\
\hline
\textbf{ی فارسی و عربی} & \lr{ي} و «ی» & «ی» \\
\hline
\textbf{ک فارسی و عربی} & \lr{ك} و «ک» & «ک» \\
\hline
\textbf{پیشوند «می»} & «می رود»، «میرود» و «می‌رود» & «می‌رود» \\
\hline
\textbf{پسوند جمع «ها»} & «کتاب ها»، «کتابها» و «کتاب‌ها» & «کتاب‌ها» \\
\hline
\textbf{فاصله‌های اضافی} & «سلام\ \ \ دنیا» و «سلام دنیا» & «سلام دنیا» \\
\hline
\textbf{اعراب و حرکات} & «مُدل» و «مدل» & «مدل» \\
\hline
\textbf{نشانه‌گذاری} & «سلام، دنیا!» و «سلام دنیا» & «سلام دنیا» \\
\hline
\end{tabular}
\endgroup
\end{table}


\زیرقسمت{کانال تلفنی و VoIP}

ناهمخوانی دامنه زمانی رخ می‌دهد که توزیع داده آموزش با ورودی در زمان استفاده متفاوت باشد. گفتار تلفنی نسبت به گفتار ضبط‌شده عمومی از نظر بداهگی، طول، مکث، نویز محیط، پاسخ میکروفن و زنجیره تخریب‌ها در کانال تفاوت دارد. حتی با وجود یک مدل زبانی قوی نیز دقت مدل بازشناسی در برابر حذف یا تغییر نشانه‌های آکوستیکی مصون نیست و مقاومت مخابراتی مسئله‌ای فراتر از شناخت واژگان فارسی است.

گفتار در یک تماس تلفنی یا \lr{VoIP} پیش از رسیدن به گیرنده از چند مرحله عبور می‌کند و در هر مرحله ممکن است بخشی از اطلاعات آکوستیکی آن تغییر کند یا از بین برود. مهم‌ترین عوامل را می‌توان به‌صورت زیر خلاصه کرد:

\شروع{فقرات}
\فقره \textbf{محدودیت پهنای‌باند:} در تلفن باریک‌باند\پاورقی{narrowband}، تنها بخشی از طیف فرکانسی گفتار منتقل می‌شود. در نتیجه، بخشی از اطلاعات فرکانسی موجود در ضبط اصلی حذف می‌شود و بازنمونه‌برداری بعدی نمی‌تواند آن را بازیابی کند.
\فقره \textbf{کدگذاری صوت:} پیش از انتقال، گفتار با یک کدک صوتی فشرده می‌شود. نوع کدک، نرخ بیت و پهنای‌باند آن بر شکل تخریب ایجادشده مؤثر است. برخی کدک‌ها شکل موج را مستقیماً کد می‌کنند و برخی دیگر گفتار را بر اساس ویژگی‌های آن فشرده می‌کنند.
\فقره \textbf{انتقال در شبکه:} در تماس‌های مبتنی بر \lr{IP}، صوت کدشده در قالب بسته‌های شبکه منتقل می‌شود. تأخیر زیاد یا از دست رفتن بسته‌ها می‌تواند باعث حذف یا تغییر بخش‌هایی از گفتار شود.
\فقره \textbf{بازسازی در گیرنده:} در صورت افت بسته، روش‌های پنهان‌سازی افت بسته\پاورقی{Packet-Loss Concealment (PLC)} تلاش می‌کنند بخش ازدست‌رفته را تخمین بزنند. این بازسازی همیشه دقیق نیست و ممکن است باعث تکرار، کشیدگی، قطع کوتاه یا تغییر شکل بخشی از گفتار شود.
\فقره \textbf{عوامل محیطی و سطح سیگنال:} نویز محیط، کیفیت میکروفن و تغییر سطح یا بهره سیگنال نیز می‌توانند همراه با عوامل بالا کیفیت صوت دریافتی را تغییر دهند.
\پایان{فقرات}

جدول~\ref{tab:telecom-codecs} کدک‌های در نظر گرفته‌شده در این پژوهش و کاربرد معمول آن‌ها را به‌طور خلاصه نشان می‌دهد.

\begin{table}[htbp]
\centering
\caption{کدک‌های مخابراتی مورد استفاده در شبیه‌سازی داده‌ها}
\label{tab:telecom-codecs}
\begingroup
\scriptsize
\setlength{\tabcolsep}{4pt}
\renewcommand{\arraystretch}{1.35}
\begin{tabular}{|p{0.15\textwidth}|p{0.20\textwidth}|p{0.27\textwidth}|p{0.28\textwidth}|}
\hline
\textbf{کدک} & \textbf{نوع/پهنای‌باند} & \textbf{کاربرد رایج} & \textbf{ویژگی کلی} \\
\hline
\lr{G.711 A-law} & باریک‌باند & شبکه‌های تلفنی ثابت و برخی سامانه‌های \lr{VoIP}، به‌ویژه در اروپا & کدگذاری \lr{PCM} با فشرده‌سازی لگاریتمی؛ نرخ بیت نسبتاً بالا و پردازش ساده \مرجع{itu1988g711} \\
\hline
\lr{G.711 $\mu$-law} & باریک‌باند & شبکه‌های تلفنی و \lr{VoIP}، به‌ویژه در آمریکای شمالی و ژاپن & مشابه \lr{A-law} با قانون کوانتیزه‌سازی متفاوت \مرجع{itu1988g711} \\
\hline
\lr{GSM Full Rate} & باریک‌باند & شبکه‌های سلولی نسل دوم \lr{GSM} & فشرده‌سازی گفتار با استفاده از پیش‌بینی و مدل‌سازی سیگنال برای کاهش نرخ بیت \مرجع{etsi2000gsm0610} \\
\hline
\lr{AMR-NB} & باریک‌باند & تماس صوتی در شبکه‌های سلولی نسل‌های دوم و سوم & کدک تطبیقی گفتار با چند نرخ بیت، طراحی‌شده برای ارتباطات سیار \مرجع{3gpp2024amrnb} \\
\hline
\lr{AMR-WB} & پهن‌باند & تماس‌های صوتی با کیفیت بالاتر و سامانه‌های \lr{HD Voice} & نسخهٔ پهن‌باند خانوادهٔ \lr{AMR} با حفظ بخش بیشتری از طیف گفتار \مرجع{3gpp2024amrwb} \\
\hline
\lr{Opus} & باریک‌باند تا تمام‌باند & \lr{VoIP}، تماس‌های اینترنتی، کنفرانس صوتی و ارتباطات بلادرنگ & کدکی انعطاف‌پذیر با پشتیبانی از نرخ بیت و پهنای‌باند متغیر، مناسب برای انتقال اینترنتی \مرجع{valin2012opus} \\
\hline
\end{tabular}
\endgroup
\end{table}

بر اساس این مراحل، در این پایان‌نامه تلاش شده است بخشی از تغییراتی که در یک ارتباط مخابراتی واقعی رخ می‌دهند به‌صورت مصنوعی روی داده‌های آموزشی اعمال شوند. هدف این شبیه‌سازی بازسازی دقیق یک شبکه یا اپراتور خاص نیست، بلکه قرار دادن مدل در معرض طیفی از شرایطی است که در تماس‌های تلفنی و \lr{VoIP} ممکن است رخ دهند. به این ترتیب، مدل در زمان آموزش تنها با گفتار تمیز روبه‌رو نمی‌شود و می‌تواند نسبت به تغییرات ناشی از پهنای‌باند، کدک، نویز و اختلالات شبکه مقاوم‌تر شود.

\قسمت{رویکردهای مقاوم‌سازی بازشناس}

\زیرقسمت{آموزش چندشرایطی}

آموزش چندشرایطی\پاورقی{Multi-Condition Training (MCT)} روشی است که در آن بازشناس با نمونه‌هایی از گفتار در شرایط آکوستیکی و مخابراتی متفاوت آموزش می‌بیند. این شرایط می‌توانند شامل نویز، محدودیت پهنای‌باند، کدک‌های مختلف و اختلالات ناشی از انتقال صوت باشند. هدف آن است که مدل به‌جای وابستگی زیاد به شرایط ضبط یا انتقال، ویژگی‌های مرتبط با محتوای گفتار را بهتر یاد بگیرد. بنابراین، تخریب مصنوعی داده در این روش صرفاً برای افزایش تعداد نمونه‌ها نیست، بلکه برای آشنا کردن مدل با شرایطی است که انتظار می‌رود در کاربرد واقعی با آن‌ها روبه‌رو شود \مرجع{li2014overview,fernandez2022augmentation}.

ترکیب گفتار اصلی و نسخه‌های تخریب‌شده امکان مقاوم‌سازی مدل در برابر شرایط مخابراتی را فراهم می‌کند، بدون آنکه عملکرد آن روی گفتار معمولی کاملاً نادیده گرفته شود. نسبت این دو نوع داده نیز اهمیت دارد؛ اگر یک کدک، نوع نویز یا شدت تخریب بیش از حد در داده آموزشی تکرار شود، مدل ممکن است بیش از اندازه با همان وضعیت سازگار شود. از سوی دیگر، حذف نمونه‌های اصلی می‌تواند عملکرد روی گفتار عادی را کاهش دهد. آموزش چندشرایطی نیز لزوماً مقاومت در برابر هر نوع شرایط دیده‌نشده را تضمین نمی‌کند؛ به همین دلیل، اثر آن باید با مقایسه یک مدل ثابت پیش و پس از افزودن داده‌های تخریب‌شده و سپس ارزیابی روی گفتار مخابراتی واقعی سنجیده شود.

\زیرقسمت{بهسازی وظیفه‌محور}

بهسازی گفتار\پاورقی{Speech Enhancement (SE)} به فرآیند برآورد یا تولید نسخه‌ای پاک‌تر از گفتار آلوده، با هدف بهبود کیفیت یا فهم‌پذیری آن گفته می‌شود. یک بهساز می‌تواند در سطح شکل موج، طیف یا ویژگی‌هایی مانند لاگ-مل عمل کند. معیارهایی مانند PESQ،‏ STOI و SI-SDR به‌ترتیب جنبه‌هایی از کیفیت ادراکی، فهم‌پذیری و وفاداری سیگنال به مرجع را می‌سنجند، اما این اهداف لزوماً با هدف بازشناسی گفتار یکسان نیستند. حذف نویز ممکن است بخشی از نشانه‌های واجی ضعیف را نیز حذف کند یا اعوجاج‌ها و آثار پردازشی تازه‌ای ایجاد کند که برای بازشناس زیان‌بار باشند. پژوهش‌های اخیر این آثار را جدا از نویز باقی‌مانده بررسی کرده و نشان داده‌اند که آموزش همسو با ASR می‌تواند نوع خطاهای ایجادشده توسط بهساز را تغییر دهد \مرجع{iwamoto2024artifacts,ochiai2024distortions}.

در بهسازی گفتارِ بازشناسی‌محور\پاورقی{Recognition-Oriented Speech Enhancement}، هدف آموزش بهساز تنها نزدیک‌کردن خروجی به گفتار پاک نیست، بلکه بهبود عملکرد بازشناسی نیز در تابع هدف لحاظ می‌شود. زیان بازسازی، خروجی بهساز را به سیگنال پاک نزدیک می‌کند؛ زیان تطبیق ویژگی، فاصله میان بازنمایی‌های بازشناس برای گفتار بهسازی‌شده و گفتار پاک را کاهش می‌دهد؛ و زیان ASR مستقیماً به حفظ اطلاعات لازم برای پیش‌بینی متن صحیح جهت می‌دهد. این زیان‌ها می‌توانند به‌صورت جداگانه یا ترکیبی و در مراحل مختلف آموزش به کار روند. در نتیجه، سود نهایی بهساز برای بازشناسی باید با معیارهایی مانند WER و CER سنجیده شود و بهبود معیارهای سیگنالی یا ادراکی به‌تنهایی تضمین‌کنندهٔ بهبود عملکرد ASR نیست \مرجع{yu2024rose}.


\زیرقسمت{حفظ نمای اصلی و همجوشی دو نما}

در خط لوله آبشاری، بازشناس فقط خروجی بهساز را دریافت می‌کند و اطلاعات حذف‌شده در بهسازی دیگر در دسترس نیست. نگه‌داشتن هم‌زمان نمای اصلی و نمای بهسازی‌شده این محدودیت را کاهش می‌دهد. نمای اصلی شواهد دست‌نخورده اما آلوده را حفظ می‌کند و نمای دوم بر بخش‌هایی تأکید دارد که بهساز آن‌ها را پاک‌تر برآورد کرده است. مطالعات جدید درباره اتصال بهسازهای ازپیش‌آموخته به بازشناس نشان داده‌اند که استفاده مستقیم از خروجی بهساز همیشه بهترین گزینه نیست و نحوه ادغام دو مشاهده بر نتیجه اثر دارد \مرجع{wang2024bridge,cui2025realbridge}.

همجوشی می‌تواند در سطح شکل موج، ویژگی‌های ورودی یا بازنمایی‌های نهان انجام شود. در همجوشی مبتنی بر توجه متقاطع\پاورقی{Cross-Attention}، یکی از جریان‌ها با مراجعه به کلیدها و مقادیر جریان دیگر، اطلاعات مرتبط را از آن استخراج می‌کند. در مقابل، یک سازوکار دروازه‌ایِ قابل‌آموزش\پاورقی{Learnable Gating} وزن یا سهم هر یک از دو بازنمایی را در خروجی ترکیبی تعیین می‌کند؛ این وزن‌ها می‌توانند در طول زمان یا میان مؤلفه‌های بازنمایی تغییر کنند. بنابراین، توجه متقاطع برای تبادل اطلاعات میان دو جریان به کار می‌رود، درحالی‌که دروازه تعیین می‌کند مدل در هر بخش تا چه اندازه به هر جریان متکی باشد. بااین‌حال، افزایش پیچیدگی معماری زمانی موجه است که در مقایسه با یک بازشناس تک‌جریانه که با همان داده‌های چندشرایطی آموزش دیده است، بهبود قابل‌توجهی ایجاد کند. معماری و تابع دقیق همجوشی پیشنهادی این پایان‌نامه در فصل چهارم معرفی می‌شود.


\قسمت{استفاده از صوت بلندِ بدون رونویسی}

\زیرقسمت{آماده‌سازی و قطعه‌بندی صوت بلند}

منابع صوتی بدون رونویسی، مانند کتاب‌های صوتی، مصاحبه‌ها و برنامه‌های گفتاری، می‌توانند حجم زیادی از گفتار واقعی و متنوع را برای آموزش بازشناس فراهم کنند. بااین‌حال، این فایل‌ها معمولاً طولانی‌اند و نمی‌توان آن‌ها را مستقیماً به‌عنوان نمونه آموزشی یا ورودی شبه‌برچسب‌گذاری استفاده کرد. وجود سکوت، موسیقی، بخش‌های غیرگفتاری و تغییرات کیفیت نیز استفاده مستقیم از آن‌ها را دشوار می‌کند. بنابراین، پیش از تولید شبه‌برچسب، صوت بلند باید به قطعه‌های کوتاه و مناسب برای بازشناس تبدیل شود. کیفیت این قطعه‌بندی و پالایش اولیه نیز بر کیفیت شبه‌برچسب‌های تولیدشده و در نتیجه بر عملکرد بهتر داده‌های بدون رونویسی در آموزش بازشناس اثر می‌گذارد \مرجع{galvez2021peoples,chen2021gigaspeech}.

یکی از روش‌های اصلی برای این کار، آشکارسازی فعالیت گفتاری\پاورقی{Voice Activity Detection (VAD)} است. VAD نواحی دارای گفتار را از سکوت و بخش‌های غیرگفتاری جدا می‌کند و می‌تواند برای تعیین مرز اولیه قطعه‌ها به کار رود. بااین‌حال، مرز آکوستیکی همیشه با مرز زبانی یکسان نیست؛ برای مثال، مکث ممکن است در میانه جمله رخ دهد و از سوی دیگر چند جمله ممکن است بدون سکوت طولانی پشت سر هم بیان شوند. بنابراین، طول قطعه و نحوه انتخاب مرزها بر میزان حفظ پیوستگی گفتار و در نهایت کیفیت شبه‌برچسب تولیدشده اثر می‌گذارند. ارزیابی اخیر روی محتوای فارسی یوتیوب نیز نشان داده است که عملکرد روش‌های قطعه‌بندی می‌تواند به نوع محتوای صوتی وابسته باشد \مرجع{huang2022e2esegmenter,iranmanesh2026segmentation}.

\زیرقسمت{شبه‌برچسب‌گذاری و پالایش متن}

شبه‌برچسب‌گذاری\پاورقی{Pseudo-Labeling} روشی برای استفاده از صوت بدون رونویسی در آموزش است. در این روش، یک مدل معلم برای هر قطعهٔ صوتی متن تولید می‌کند و نمونه‌هایی که معیارهای پذیرش را برآورده کنند، به‌عنوان داده آموزشی برای مدل دانش‌آموز استفاده می‌شوند. مزیت اصلی این روش، افزایش حجم داده قابل استفاده بدون نیاز به رونویسی دستی است. بااین‌حال، خطاهای مدل معلم، مانند حذف، درج، جانشینی یا تکرار واژه‌ها، می‌توانند وارد برچسب‌های آموزشی شوند. بنابراین، اثربخش بودن شبه‌برچسب‌گذاری تنها به حجم داده خام وابسته نیست و کیفیت مدل معلم و روش انتخاب نمونه‌های قابل اعتماد نیز نقش مهمی دارند. پژوهش‌های خودآموزی و شبه‌برچسب‌گذاری، به‌ویژه در زبان‌های کم‌منبع، بر اهمیت پالایش نمونه‌ها و تولید بازتولیدپذیر شبه‌برچسب‌ها تأکید کرده‌اند \مرجع{kahn2020selftraining,likhomanenko2021slimipl,bhogale2024pseudolabel}.

پس از برچسب‌گذاری، نمونه‌ها را می‌توان بر اساس نشانه‌های صوتی و متنی پالایش کرد. مدت قطعه، نسبت گفتار به سکوت، زبان تشخیص‌داده‌شده، وجود تکرار، نسبت طول متن به مدت صوت و اطمینان مدل از جمله این معیارها هستند. پالایش سخت‌گیرانه می‌تواند تعداد برچسب‌های نادرست را کاهش دهد، اما در مقابل ممکن است بخشی از داده‌های دشوار یا متنوع را نیز حذف کند. در نتیجه، تعداد ساعت باقی‌مانده پس از پالایش به‌تنهایی نشان‌دهنده کیفیت پیکره نیست و باید در کنار کیفیت شبه‌برچسب‌ها در نظر گرفته شود. نگه‌داری متن اولیه تولیدشده توسط معلم و دلیل پذیرش یا رد هر قطعه نیز امکان بررسی و بازتولید این فرایند را فراهم می‌کند.

پس از پالایش، متن شبه‌برچسب‌ها می‌تواند نرمال‌سازی شود تا تفاوت‌های سطح نوشتاری که برای آموزش بازشناس اهمیتی ندارند کاهش یابند. این مرحله شامل تبدیل‌های قطعی و ازپیش‌تعریف‌شده، مانند یکسان‌سازی نویسه‌های یونی‌کد، فاصله‌ها، اعداد و نشانه‌گذاری است و نباید محتوای زبانی جمله را تغییر دهد. اصلاح زبانی وابسته به بافت، مانند تغییر واژه یا بازنویسی عبارت، مسئله‌ای جداگانه است؛ زیرا مدلی که به صوت دسترسی ندارد ممکن است متنی روان‌تر اما ناسازگار با گفتار واقعی تولید کند. به همین دلیل، در صورت استفاده از چنین اصلاحی، متن خام معلم، متن اصلاح‌شده و میزان تغییر میان آن‌ها باید جداگانه نگه‌داری شوند. پژوهش‌های اخیر در اصلاح خروجی ASR نیز بر تمایز میان روانی متن و وفاداری آن به شواهد صوتی تأکید کرده‌اند \مرجع{hu2024listen,udagawa2024conservative}.


\قسمت{جمع‌بندی}

در این فصل، مفاهیم لازم برای بررسی مسئلهٔ بازشناسی گفتار فارسی در شرایط مخابراتی معرفی شدند. ابتدا نحوهٔ کار بازشناس انتهابه‌انتها و معیارهای \lr{WER} و \lr{CER} توضیح داده شد و سپس عواملی که باعث تفاوت گفتار تلفنی و \lr{VoIP} با داده‌های معمول آموزشی، از جمله محدودیت پهنای‌باند، کدگذاری صوت، نویز و اختلالات شبکه، بررسی شدند. در ادامه، آموزش چندشرایطی و استفاده هم‌زمان از نمای اصلی و بهسازی‌شده به‌عنوان دو رویکرد برای کاهش اثر این تفاوت‌ها مطرح شدند. همچنین، فرایند تبدیل صوت بلند و بدون رونویسی به داده قابل استفاده برای آموزش، از قطعه‌بندی تا شبه‌برچسب‌گذاری و پالایش متن، توضیح داده شد.

این مفاهیم مبنای روش‌هایی هستند که در ادامه پایان‌نامه بررسی می‌شوند. در فصل بعد، پژوهش‌های پیشین مرتبط با مقاوم‌سازی بازشناس و استفاده از داده‌های بدون رونویسی مرور می‌شوند و سپس در فصل راهکار پیشنهادی، نحوه به‌کارگیری این روش‌ها برای داده‌ها و مدل‌های مورد استفاده در این پژوهش تشریح خواهد شد.
